
\chapter{Aim of the thesis}
\label{sec:aim_thesis}

This thesis was carried out within the VIBRA research group, which poses as its primary aim to develop a next-generation microscopy system capable of inspecting in two dimensions the chemical content of a biological specimen using Raman scattering techniques, more specifically to identify tumors in human biopsies with accuracy and reproducibility.

The aim pursued in this thesis is to apply machine learning methods, specifically \textit{Principal Component Analysis} (PCA) and an \textit{Artificial Neural Network} (ANN), using Python tools, on a dataset consisting of hyperspectral Raman maps acquired from
patient-derived tumor biopsies of head and neck squamous cell carcinoma (HNSCC), already prepared for computation. Through these methods, the composition of the available data is investigated, with particular interest in evaluating a binary classification pipeline to identify a defined tumor class among different tissue types.

In detail, the work will focus on the following steps:
\begin{itemize}
    \item the available data are examined both quantitatively and qualitatively, exploring their structure and statistical distribution. The biological relevance of the tissue classes present is assessed in order to define the binary classification problem;

    \item  \textit{Principal Component Analysis} is computed to project the dataset into a new subspace and capture the variance along the spectral bands, to assess whether a reduced representation is appropriate. Moreover, this could enable the identification of important features that can be used for data separation and inspection;

    \item three different \textit{Artificial Neural Network} architectures are implemented, and their performance is evaluated by testing them on two distinct binary classification tasks, both of which classify data into tumor and non-tumor classes. The aim is to determine which of the tested architectures performs best in terms of prediction results and computational complexity.
\end{itemize}

